% ---------------------------------------------------------------------------
% Section: Two dimensions
% ---------------------------------------------------------------------------
% Section-local macros (unique names, \providecommand only).
\providecommand{\TLat}{\mathcal{L}}          % finite graph (V,E)
\providecommand{\TCyc}{\mathcal{Z}}          % cycle space
\providecommand{\TFlux}{\Phi}                % flux configurations
\providecommand{\TGauge}{\Theta}             % gauge group {+-1}^V
\providecommand{\Tpi}{\pi}                   % fermion parity label
\providecommand{\Tclus}{\mathcal{C}}         % equal-charge cluster

\section{Two dimensions}\label{sec:2d}

The model of \cite{SKKM2018} is defined on an arbitrary graph.  In this section we prove the exact
free-fermion mapping on every finite connected graph, with the complete list of superselection labels
(charges, fluxes through plaquettes and, on a torus, the non-contractible Wilson loops) and the parity
constraint (\cref{thm:2dmap}).  We then determine the sector laws of two families of product initial
states (\cref{prop:2dlaws}): $\sigma^z$-polarized links give the two-dimensional Anderson model with
Bernoulli potential, and $\sigma^x$-polarized links with a checkerboard charge-density wave give binary
random $\pi$-flux with a staggered potential.  At strong coupling we prove a cluster-confinement estimate,
valid up to times of order $h/J^2$ (\cref{prop:cluster}), and we explain why it cannot be pushed further.
Finally we summarize what is (not) known about the two-dimensional one-body problems.  In contrast to one
dimension, we do \emph{not} obtain sector-averaged localization for any $h$ in two dimensions: the
required one-body input is not available (\cref{sec:2d:status}).

\subsection{The model on a graph and the exact mapping}\label{sec:2d:model}

\paragraph{Setting.}
Let $\TLat=(V,E)$ be a finite connected graph without loops; multiple edges are allowed.  Every edge
$b\in E$ has two distinct endpoints, $\partial b=\{i,j\}$, and we write $b\ni i$ if $i\in\partial b$.
The sites carry spinless fermions $f_i$, $i\in V$, with Fock space $\cF_V$, $n_i=f_i^\dagger f_i$ and
$N_f=\sum_in_i$; after fixing an enumeration of $V$, the occupation basis $\ket n$, $n\in\{0,1\}^V$, is
defined by the ordering convention \eqref{eq:occbasis}, and we write $\abs n:=\sum_in_i$.  We denote by
$\cF_V^{(\Tpi)}$, $\Tpi=\pm1$, the eigenspaces of the fermion parity $(-1)^{N_f}$.  The edges carry spins
$\frac12$; $\ket s_z$, $s\in\{\pm1\}^E$, is the $\sigma^z$-basis, and $\ket{n,s}:=\ket n\otimes\ket s_z$.
For $J>0$ and $h\in\R\setminus\{0\}$ let
\begin{equation}\label{eq:H2d}
  H_\TLat:=-J\sum_{b\in E,\ \partial b=\{i,j\}}\sigma^z_b\bigl(f_i^\dagger f_j+f_j^\dagger f_i\bigr)
  -h\sum_{i\in V}A_i,\qquad A_i:=\prod_{b\ni i}\sigma^x_b ,
\end{equation}
and define the Gauss operators and Wilson loops
\begin{equation}\label{eq:G2d}
  G_i:=(-1)^{n_i}A_i\quad(i\in V),\qquad W_\gamma:=\prod_{b\in\gamma}\sigma^z_b\quad(\gamma\in\TCyc(\TLat)).
\end{equation}
Here $\TCyc(\TLat)$ is the \emph{cycle space}: the set of edge sets $\gamma\subset E$ such that every
vertex is an endpoint of an even number of edges of $\gamma$.  It is a vector space over $\mathbb F_2$
under symmetric difference.  \eqref{eq:H2d} is \cite[Eqs.~(1)--(2)]{SKKM2018} with $J_{jk}\equiv J$ and
$h_j\equiv h$ (their sum over $\langle jk\rangle$ contains both orientations); $G_i$ is their conserved
charge $\hat q_i=(-1)^{\hat n_i}\hat A_i$, and $W_\gamma$ is their plaquette operator $\hat B_p$ when $\gamma$
is a plaquette and their Wilson loop $\hat\Gamma_n$ (their Eq.~(6)) when $\gamma$ winds around a torus.  For
the square lattice with open boundaries the $A_i$ of boundary sites are the ``incomplete stars'' of
\cite[Fig.~1]{SKKM2018}.  On a path with $N$ sites and $N-1$ links, \eqref{eq:H2d} is the ``open'' chain
of \cref{rem:boundary}(d); the boundary link of \cref{sec:model} is a one-dimensional device that we do
not use here.

\paragraph{Fluxes, gauge transformations, one-body operators.}
Let $\beta:=\abs E-\abs V+1$.  A \emph{flux configuration} is a group homomorphism
$\varphi:\TCyc(\TLat)\to\{\pm1\}$; they form a group $\TFlux(\TLat)$.  A sign configuration
$s\in\{\pm1\}^E$ has the flux $\varphi_s(\gamma):=\prod_{b\in\gamma}s_b$.  The gauge group
$\TGauge:=\{\pm1\}^V$ acts on sign configurations by $(\theta\cdot s)_b:=\theta_i\theta_js_b$ for
$\partial b=\{i,j\}$.  For $q\in\{\pm1\}^V$ and $s\in\{\pm1\}^E$ we define on $\ell^2(V)$
\begin{equation}\label{eq:hqs}
  K_s:=-J\sum_{b\in E,\ \partial b=\{i,j\}}s_b\bigl(\ket i\bra j+\ket j\bra i\bigr),\qquad
  h_{q,s}:=K_s+2h\,\diag(q),
\end{equation}
the hopping operator with binary $\pi$-fluxes $\varphi_s$ and the binary potential $\pm2h$.  With
$D_\theta:=\diag(\theta)$ one has $D_\theta h_{q,s}D_\theta=h_{q,\theta\cdot s}$, so by
\cref{lem:2dgauge}(i) the spectrum of $h_{q,s}$ and the moduli $\abs{f(h_{q,s})_{xy}}$ depend on $s$ only
through $\varphi_s$.  As in \cref{sec:model}, $\dGamma(a):=\sum_{x,y\in V}a_{xy}f_x^\dagger f_y$, and
$\Gamma(u)$ is the second quantization of a unitary $u$.

\begin{lemma}[Fluxes and gauge orbits]\label{lem:2dgauge}
\leavevmode
\begin{enumerate}[label=(\roman*)]
\item For $s,s'\in\{\pm1\}^E$: $\varphi_s=\varphi_{s'}$ if and only if $s'=\theta\cdot s$ for some
$\theta\in\TGauge$, and $\theta$ is then unique up to a global sign.
\item Let $T\subset E$ be a spanning tree.  For $b\in E\setminus T$ let $\gamma_b$ be the edge set
consisting of $b$ and the path in $T$ joining the endpoints of $b$.  Then
$\{\gamma_b:b\in E\setminus T\}$ is a basis of $\TCyc(\TLat)$; in particular
$\dim_{\mathbb F_2}\TCyc(\TLat)=\beta$ and $\abs{\TFlux(\TLat)}=2^\beta$.
\item The map $s\mapsto\varphi_s$ is onto $\TFlux(\TLat)$, and each fibre
$O_\varphi:=\{s:\varphi_s=\varphi\}$ is a $\TGauge$-orbit with $2^{\abs V-1}$ elements.
\end{enumerate}
\end{lemma}

\begin{proof}
(ii) Each $\gamma_b$ is a cycle (every vertex of a simple closed path has degree $2$), and $\gamma_b$ is the
only one of these sets containing the non-tree edge $b$; hence they are linearly independent.  Given
$\gamma\in\TCyc$, the set $\gamma':=\gamma\,\triangle\,\triangle_{b\in\gamma\setminus T}\gamma_b$ is a cycle
contained in $T$.  A non-empty edge set in a tree has a vertex of degree one (an endpoint of a longest
path in it), so $\gamma'=\emptyset$, and $\gamma=\triangle_{b\in\gamma\setminus T}\gamma_b$.  Since
$\abs T=\abs V-1$, $\dim\TCyc=\abs{E\setminus T}=\beta$; a homomorphism to $\{\pm1\}$ is determined by
its arbitrary values on a basis, so $\abs{\TFlux}=2^\beta$.

(i) If $s'=\theta\cdot s$ and $\gamma\in\TCyc$, then
$\varphi_{s'}(\gamma)=\varphi_s(\gamma)\prod_{i\in V}\theta_i^{d_\gamma(i)}=\varphi_s(\gamma)$, where the
degree $d_\gamma(i)$ of $i$ in $\gamma$ is even.  Conversely let $\varphi_s=\varphi_{s'}$ and $r:=ss'$
(pointwise), so that $\varphi_r\equiv1$.  Fix a root $v_0$ and define $\theta_v:=\prod_{b\in T(v_0,v)}r_b$,
the product over the tree path from $v_0$ to $v$.  For an edge $b$ with $\partial b=\{i,j\}$, if $b\in T$
then $T(v_0,i)\triangle T(v_0,j)=\{b\}$; if $b\notin T$ then $T(v_0,i)\triangle T(v_0,j)\triangle\{b\}=\gamma_b$.
In both cases $r_b=\prod_{T(v_0,i)\triangle T(v_0,j)}r=\theta_i\theta_j$ (use $r^2=1$ and
$\varphi_r(\gamma_b)=1$), i.e.\ $s'=\theta\cdot s$.  If $\theta\cdot s=\theta'\cdot s$, then
$(\theta\theta')_i=(\theta\theta')_j$ for every edge, so $\theta\theta'$ is constant because $\TLat$ is
connected.

(iii) Given $\varphi$, set $s_b:=1$ for $b\in T$ and $s_b:=\varphi(\gamma_b)$ for $b\notin T$; then
$\varphi_s=\varphi$ on the basis of (ii), hence everywhere.  By (i) the fibres are orbits, and the
stabilizer of every $s$ is $\{\pm(1,\dots,1)\}$, so each orbit has $2^{\abs V}/2$ elements.
\end{proof}

\begin{theorem}[Exact mapping on graphs; planar boxes and tori]\label{thm:2dmap}
Let $\TLat=(V,E)$ be a finite connected graph without loops.
\begin{enumerate}[label=(\alph*)]
\item The $G_i$ ($i\in V$) and $W_\gamma$ ($\gamma\in\TCyc(\TLat)$) are self-adjoint unitaries; they commute
with each other, with $N_f$ and with $H_\TLat$; and
\begin{equation}\label{eq:2dparity}
  \prod_{i\in V}G_i=(-1)^{N_f}.
\end{equation}
\item For $q\in\{\pm1\}^V$ and $\varphi\in\TFlux(\TLat)$ let $P_{q,\varphi}$ be the joint spectral projection
onto $\{G_i=q_i\ \forall i,\ W_\gamma=\varphi(\gamma)\ \forall\gamma\}$, and let $\Tpi_q:=\prod_iq_i$.  Then
$\sum_{q,\varphi}P_{q,\varphi}=\one$, $\operatorname{rank}P_{q,\varphi}=2^{\abs V-1}$ for \emph{every}
$(q,\varphi)$, and $(-1)^{N_f}P_{q,\varphi}=\Tpi_qP_{q,\varphi}$.  For every $s$ with $\varphi_s=\varphi$ there
is a unitary $V_{q,s}:\cF_V^{(\Tpi_q)}\to\operatorname{Ran}P_{q,\varphi}$ which maps each $\ket n$ to a joint
eigenvector of the $n_i$ with the same eigenvalues, and such that, on $\cF_V^{(\Tpi_q)}$,
\begin{equation}\label{eq:2ddict}
  V_{q,s}^*\,\sigma^z_bf_i^\dagger f_j\,V_{q,s}=s_b\,f_i^\dagger f_j\quad(\partial b=\{i,j\}),\qquad
  V_{q,s}^*A_iV_{q,s}=q_i(-1)^{n_i},
\end{equation}
and consequently
\begin{equation}\label{eq:2dsector}
  V_{q,s}^*\,H_\TLat\,V_{q,s}=\Bigl(\dGamma(h_{q,s})-h\sum_{i\in V}q_i\Bigr)\Big|_{\cF_V^{(\Tpi_q)}} .
\end{equation}
For $\theta\in\TGauge$, $V_{q,\theta\cdot s}=q(\theta)\,V_{q,s}\,\Gamma(D_\theta)$ with
$q(\theta):=\prod_{i:\theta_i=-1}q_i$.
\item \emph{Planar box.}  Let $\TLat$ be the box $\{1,\dots,L_1\}\times\{1,\dots,L_2\}$ with nearest-neighbour
edges, and let the plaquettes $p$ be its unit squares, with boundaries $\partial p$ and $W_p:=W_{\partial p}$.
Then $\beta=(L_1-1)(L_2-1)$, and $\{\partial p\}_p$ is a basis of $\TCyc(\TLat)$.  Hence the sectors are
labelled by $(q,(W_p)_p)\in\{\pm1\}^V\times\{\pm1\}^{\#\text{plaquettes}}$, all $2^{\abs V+\beta}$ labels
occur, and in the sector $(q,(W_p))$ the matter consists of free fermions with the binary potential $2hq_i$
and the binary $\pi$-flux $W_p=\prod_{b\in\partial p}s_b\in\{\pm1\}$ through $p$, restricted to the parity
sector $(-1)^{N_f}=\prod_iq_i$.
\item \emph{Torus.}  Let $\TLat$ be the torus $(\Z/L_1\Z)\times(\Z/L_2\Z)$, $L_1,L_2\ge3$, with
nearest-neighbour edges, plaquettes $p$ as in (c), and let $\gamma_1$ (resp.\ $\gamma_2$) be the set of
horizontal edges of the row $\{y=0\}$ (resp.\ the vertical edges of the column $\{x=0\}$).  Then
$\beta=L_1L_2+1$, $\prod_pW_p=\one$, and $\{\partial p:p\neq p_0\}\cup\{\gamma_1,\gamma_2\}$ is a basis of
$\TCyc(\TLat)$ for any fixed plaquette $p_0$.  Hence the sectors are labelled by $q$, the plaquette fluxes
$(W_p)_p$ subject to $\prod_pW_p=1$, and the two non-contractible Wilson loops $W_{\gamma_1},W_{\gamma_2}$,
which are independent labels; all $2^{\abs V+\beta}$ labels occur, and in $h_{q,s}$ the non-contractible
loops act as $\Z_2$ twists of the boundary conditions in the two directions.
\end{enumerate}
\end{theorem}

\begin{proof}
For $\theta\in\TGauge$ let $V_\theta:=\{i:\theta_i=-1\}$, $q(\theta):=\prod_{i\in V_\theta}q_i$,
$\varepsilon_n(\theta):=\prod_{i\in V_\theta}(-1)^{n_i}$ and $\mathcal G_\theta:=\prod_{i\in V_\theta}G_i$.
Since $\sigma^x$ flips $\sigma^z$-eigenstates and $\TLat$ has no loops, $A_i\ket{n,s}=\ket{n,\theta^{(i)}\cdot s}$,
where $\theta^{(i)}$ equals $-1$ at $i$ only.  Hence
\begin{equation}\label{eq:Gtheta}
  \mathcal G_\theta\ket{n,s}=\varepsilon_n(\theta)\,\ket{n,\theta\cdot s}\qquad(\theta\in\TGauge).
\end{equation}

(a) $G_i$ and $W_\gamma$ are products of commuting self-adjoint unitaries.  The $G_i$ are built from the
mutually commuting operators $\sigma^x_b$ and $(-1)^{n_k}$; the $W_\gamma$ are diagonal in the basis
$\ket{n,s}$, so they commute with each other, with $N_f$ and with every $(-1)^{n_k}$; and
$A_iW_\gamma=(-1)^{d_\gamma(i)}W_\gamma A_i=W_\gamma A_i$ because the degree $d_\gamma(i)$ is even.  The
electric term commutes with all $\sigma^x_b$, $(-1)^{n_k}$ and $W_\gamma$.  Let
$T_b:=\sigma^z_b(f_i^\dagger f_j+f_j^\dagger f_i)$, $\partial b=\{i,j\}$.  It commutes with every
$W_\gamma$.  It anticommutes with $A_k$ iff $b\ni k$, and with $(-1)^{n_k}$ iff $k\in\{i,j\}$ (because
$(-1)^{n_i}f_i^\dagger=-f_i^\dagger(-1)^{n_i}$ and $(-1)^{n_i}$ commutes with $f_j$, $j\ne i$); so it
commutes with each $G_k=(-1)^{n_k}A_k$.  Everything commutes with $N_f$.  Finally
$\prod_iG_i=(-1)^{N_f}\prod_iA_i=(-1)^{N_f}\prod_{b\in E}(\sigma^x_b)^2=(-1)^{N_f}$, since every edge has
exactly two endpoints.

(b) \emph{Step 1: joint eigenvectors.}  Let $P_q:=\prod_i\frac12(\one+q_iG_i)=2^{-\abs V}\sum_{\theta\in\TGauge}q(\theta)\mathcal G_\theta$.
By \eqref{eq:Gtheta}, $P_q\ket{n,s}=2^{-\abs V}\sum_\theta q(\theta)\varepsilon_n(\theta)\ket{n,\theta\cdot s}$.  The
elements $\theta$ and $-\theta$ give the same vector $\ket{n,\theta\cdot s}$, and
$q(-\theta)\varepsilon_n(-\theta)=\Tpi_q(-1)^{\abs n}q(\theta)\varepsilon_n(\theta)$.  Hence $P_q\ket{n,s}=0$ if
$(-1)^{\abs n}\neq\Tpi_q$, while for $(-1)^{\abs n}=\Tpi_q$
\begin{equation}\label{eq:Psi2d}
  \Psi_{n,q,s}:=2^{(\abs V-1)/2}P_q\ket{n,s}
  =2^{-(\abs V-1)/2}\sum_{[\theta]\in\TGauge/\{\pm1\}}q(\theta)\varepsilon_n(\theta)\,\ket{n,\theta\cdot s},
\end{equation}
where the summand depends only on the class $[\theta]=\{\theta,-\theta\}$.  By \cref{lem:2dgauge}(i),(iii),
$[\theta]\mapsto\theta\cdot s$ is a bijection onto the orbit $O_{\varphi_s}$, so \eqref{eq:Psi2d} is a normalized
sum of $2^{\abs V-1}$ orthonormal vectors with unimodular coefficients: $\Psi_{n,q,s}$ is a unit vector.  It
satisfies $G_i\Psi_{n,q,s}=q_i\Psi_{n,q,s}$ (as $G_iP_q=q_iP_q$), $W_\gamma\Psi_{n,q,s}=\varphi_s(\gamma)\Psi_{n,q,s}$
(every $\ket{n,\theta\cdot s}$ has flux $\varphi_s$), and it is an eigenvector of every occupation operator, with
the same eigenvalues as $\ket n$.

\emph{Step 2: orthonormal bases.}  Fix $\varphi$ and $s\in O_\varphi$.  Since the $W_\gamma$ are diagonal in the
basis $\ket{n,s'}$, their joint eigenspace with eigenvalues $\varphi$ is spanned by $\{\ket{n,s'}:s'\in O_\varphi\}$.
The joint eigenvalues of the commuting self-adjoint unitaries $(W_\gamma)_\gamma$ form a homomorphism
$\gamma\mapsto\varphi(\gamma)$, because $W_{\gamma\triangle\gamma'}=W_\gamma W_{\gamma'}$; hence
$\sum_{q,\varphi}P_{q,\varphi}=\one$.  As $P_q$ commutes with the $W_\gamma$, $\operatorname{Ran}P_{q,\varphi}$ is
spanned by the vectors $P_q\ket{n,\theta\cdot s}$, $\theta\in\TGauge$.  By \eqref{eq:Gtheta},
$\ket{n,\theta\cdot s}=\varepsilon_n(\theta)\mathcal G_\theta\ket{n,s}$, and $P_q\mathcal G_\theta=q(\theta)P_q$, so
\begin{equation}\label{eq:Psigauge}
  P_q\ket{n,\theta\cdot s}=q(\theta)\varepsilon_n(\theta)\,P_q\ket{n,s}.
\end{equation}
Therefore $\{\Psi_{n,q,s}:(-1)^{\abs n}=\Tpi_q\}$ spans $\operatorname{Ran}P_{q,\varphi}$; it is orthonormal because
distinct $n$ give orthogonal vectors.  This proves $\operatorname{rank}P_{q,\varphi}=2^{\abs V-1}$ and
$(-1)^{N_f}=\Tpi_q$ on $\operatorname{Ran}P_{q,\varphi}$, and $V_{q,s}\ket n:=\Psi_{n,q,s}$ defines a unitary
$\cF_V^{(\Tpi_q)}\to\operatorname{Ran}P_{q,\varphi}$.  By \eqref{eq:Psigauge},
$\Psi_{n,q,\theta\cdot s}=q(\theta)\varepsilon_n(\theta)\Psi_{n,q,s}$, which is the stated change of representative,
since $\Gamma(D_\theta)\ket n=\varepsilon_n(\theta)\ket n$.

\emph{Step 3: the Hamiltonian.}  Since $A_i=(-1)^{n_i}G_i$, $A_i\Psi_{n,q,s}=q_i(-1)^{n_i}\Psi_{n,q,s}$.  For the
hopping let $\partial b=\{i,j\}$ and $n_j=1$, $n_i=0$, so that $f_i^\dagger f_j\ket n=\eta\ket{n'}$ with
$n':=n-e_j+e_i$ and a sign $\eta$ given by \eqref{eq:fsign}; if $n_j=0$ or $n_i=1$ then $f_i^\dagger f_j\ket n=0$.
As $n'$ and $n$ differ exactly at $i$ and $j$, $\varepsilon_{n'}(\theta)=\theta_i\theta_j\varepsilon_n(\theta)$, and
$\sigma^z_b\ket{\theta\cdot s}_z=\theta_i\theta_js_b\ket{\theta\cdot s}_z$.  Hence, by \eqref{eq:Psi2d},
\[
  \sigma^z_bf_i^\dagger f_j\,\Psi_{n,q,s}
  =2^{-(\abs V-1)/2}\sum_{[\theta]}q(\theta)\varepsilon_n(\theta)\theta_i\theta_js_b\,\eta\,\ket{n',\theta\cdot s}
  =s_b\,\eta\,\Psi_{n',q,s},
\]
and both sides vanish in the remaining cases.  This is \eqref{eq:2ddict}.  Summing,
$V_{q,s}^*H_\TLat V_{q,s}=-J\sum_bs_b(f_i^\dagger f_j+f_j^\dagger f_i)-h\sum_iq_i(1-2n_i)$, which is
\eqref{eq:2dsector}.

(c) The box is connected, $\abs V=L_1L_2$ and $\abs E=(L_1-1)L_2+L_1(L_2-1)$, so $\beta=(L_1-1)(L_2-1)$, the
number of plaquettes.  Each $\partial p$ is a cycle.  If $\triangle_{p\in P'}\partial p=\emptyset$ for a non-empty set
$P'$ of plaquettes, pick $p\in P'$ whose lower-left corner has maximal first coordinate; the right edge of $p$
lies in $\partial p$ and otherwise only in the boundary of the plaquette to the right of $p$, which is not in
$P'$, so it survives in the symmetric difference, a contradiction.  So the $\partial p$ are $\beta$ independent
vectors, hence a basis, and $\varphi\mapsto(\varphi(\partial p))_p$ is a bijection
$\TFlux\to\{\pm1\}^{\#\text{plaquettes}}$.  The rest follows from (b).

(d) For $L_1,L_2\ge3$ the torus is a simple $4$-regular connected graph with $\abs V=L_1L_2$ and $\abs E=2L_1L_2$,
so $\beta=L_1L_2+1$.  Every edge lies in exactly two plaquettes, so $\triangle_p\partial p=\emptyset$, i.e.\
$\prod_pW_p=\one$.  If $\triangle_{p\in P'}\partial p=\emptyset$ with $\emptyset\neq P'\neq$ all plaquettes, then,
the dual torus being connected, some edge lies in a plaquette of $P'$ and in one outside $P'$, and survives:
a contradiction.  Hence the $\partial p$, $p\ne p_0$, are independent and span the same space $B$ as all
$\partial p$.  Let $C_1$ be the set of horizontal edges joining $(0,y)$ and $(1,y)$, $y\in\Z/L_2\Z$, and $C_2$
the set of vertical edges joining $(x,0)$ and $(x,1)$.  The linear forms $\ell_k(\gamma):=\abs{\gamma\cap C_k}\bmod2$
vanish on every $\partial p$ (a plaquette contains $0$ or $2$ edges of $C_k$), and
$\ell_1(\gamma_1)=\ell_2(\gamma_2)=1$, $\ell_1(\gamma_2)=\ell_2(\gamma_1)=0$.  Hence $\gamma_1,\gamma_2$ are independent
modulo $B$, and $\{\partial p:p\ne p_0\}\cup\{\gamma_1,\gamma_2\}$ consists of $L_1L_2+1=\beta$ independent cycles, a
basis.  Choosing $s$ as in \cref{lem:2dgauge}(iii), $W_{\gamma_1}$ and $W_{\gamma_2}$ are the products of the hopping
signs around the two non-contractible directions, i.e.\ twisted boundary conditions.
\end{proof}

\begin{remark}[Comparison with \cite{SKKM2018}]\label{rem:2dSKKM}
SKKM introduce dual spins by $\hat\tau^z_j=\hat A_j$, $\hat\tau^x_j\hat\tau^x_k=\hat\sigma^z_{jk}$
\cite[Eq.~(5)]{SKKM2018}, which forces $\hat B_p=1$ and $\hat\Gamma_n=1$; other Wilson-loop sectors are reached by
modifying the duality along reference lines \cite[Eq.~(19)]{SKKM2018}, and the global constraint
$\prod_i\hat q_i=(-1)^{\hat N_f}$ is their Eq.~(7).  \Cref{thm:2dmap} treats all sectors at once: the dimension
count $2^{\abs V}\cdot2^\beta\cdot2^{\abs V-1}=2^{\abs V+\abs E}$ shows that the labels $(q,\varphi)$ exhaust
$\cH$, every sector $(q,\varphi)$ carries the free-fermion Hamiltonian $\dGamma(h_{q,s})$, $s\in O_\varphi$, on the
parity sector $\Tpi_q$, and the only constraint is the parity relation \eqref{eq:2dparity}.  Unlike in the chain with a boundary link
(\cref{lem:gauss}), the charges of a graph without ``dangling'' links can never be independent of the matter
parity.
\end{remark}

\subsection{Sector laws of product initial states}\label{sec:2d:laws}

For an operator $O$ on $\cH$ commuting with all $G_i$ and all $W_\gamma$ we write
$O_{q,s}:=V_{q,s}^*OV_{q,s}$, an operator on $\cF_V^{(\Tpi_q)}$; for every function $F$ of the occupation
numbers, $F_{q,s}=F$.  We write $\Psi(t):=\e^{-\ii tH_\TLat}\Psi_0$.

\begin{proposition}[Sector laws]\label{prop:2dlaws}
Let $\TLat$ be as in \cref{thm:2dmap}.
\begin{enumerate}[label=(\alph*)]
\item \emph{$\sigma^z$-product links: the Anderson--Bernoulli model.}  Let $\Psi_0=\psi\otimes\ket{s}_z$ with
$s\in\{\pm1\}^E$ and $\psi\in\cF_V^{(\Tpi)}$ of definite parity $\Tpi$.  Then $P_{q,\varphi}\Psi_0=0$ unless
$\varphi=\varphi_s$ and $\Tpi_q=\Tpi$, and for these sectors $P_{q,\varphi_s}\Psi_0=2^{-(\abs V-1)/2}\,V_{q,s}\psi$.  So the
flux is deterministic, the sector law of the charges is uniform on the parity hyperplane $\{q:\prod_iq_i=\Tpi\}$,
every sector carries the matter state $\psi$, and for every $O$ commuting with all $G_i,W_\gamma$
\begin{equation}\label{eq:2dlawz}
  \langle\Psi(t),O\Psi(t)\rangle=2^{1-\abs V}\sum_{q:\ \prod_iq_i=\Tpi}
  \bigl\langle\Gamma(\e^{-\ii th_{q,s}})\psi,\,O_{q,s}\,\Gamma(\e^{-\ii th_{q,s}})\psi\bigr\rangle .
\end{equation}
The law of $q$ has density at most $2$ with respect to the i.i.d.\ Bernoulli$(1/2)$ law.  For the
$\sigma^z$-polarized state $\ket{\uparrow\cdots\uparrow}_z$, i.e.\ $s=+:=(1,\dots,1)$, all fluxes vanish and the
one-body operators are the Anderson--Bernoulli operators $h_{q,+}=-J\Delta+2h\diag(q)$, with $\Delta$ the adjacency
operator of $\TLat$.
\item \emph{$\sigma^x$-polarized links: random $\pi$-flux.}  Let $\Psi_0=\ket{n^0}\otimes\ket{+\cdots+}$, where
$\sigma^x_b\ket{+\cdots+}=\ket{+\cdots+}$ for all $b$.  Let $q^0_i:=(-1)^{n^0_i}$.  Then $P_{q,\varphi}\Psi_0=0$
unless $q=q^0$, and for every $\varphi\in\TFlux(\TLat)$ and every $s\in O_\varphi$,
$P_{q^0,\varphi}\Psi_0=2^{-\beta/2}\,V_{q^0,s}\ket{n^0}$.  So the charges are deterministic, the flux
configuration is uniformly distributed on $\TFlux(\TLat)$ (on a box: i.i.d.\ uniform plaquette fluxes; on a
torus: $(W_p)$ uniform on $\{\prod_pW_p=1\}$, and $W_{\gamma_1},W_{\gamma_2}$ independent and uniform), every sector
carries the matter state $\ket{n^0}$, and
\begin{equation}\label{eq:2dlawx}
  \langle\Psi(t),O\Psi(t)\rangle=2^{-\beta}\sum_{\varphi\in\TFlux}
  \bigl\langle\Gamma(\e^{-\ii th_{q^0,s^\varphi}})n^0,\,O_{q^0,s^\varphi}\,\Gamma(\e^{-\ii th_{q^0,s^\varphi}})n^0\bigr\rangle ,
  \qquad s^\varphi\in O_\varphi .
\end{equation}
\item \emph{Checkerboard.}  In (b) let $\TLat$ be bipartite with sublattice signs $\epsilon_i\in\{\pm1\}$
($\epsilon_i=-\epsilon_j$ whenever $\partial b=\{i,j\}$), and $n^0_i=\frac12(1+\epsilon_i)$.  Let
$\Gamma_\epsilon:=\diag(\epsilon)$.  Then $h_{q^0,s}=K_s-2h\Gamma_\epsilon$ with
$K_s\Gamma_\epsilon=-\Gamma_\epsilon K_s$, and for every $s$:
\begin{enumerate}[label=(\roman*)]
\item $h_{q^0,s}^2=K_s^2+4h^2$;
\item $\sigma(h_{q^0,s})\subset(-\infty,-2\abs h]\cup[2\abs h,\infty)$, and for every Borel set $B\subset[0,\infty)$,
$\one_B(\abs{h_{q^0,s}})=\one_{g^{-1}(B)}(\abs{K_s})$ with $g(e):=\sqrt{e^2+4h^2}$.  In particular
$\one_{\{2\abs h\}}(\abs{h_{q^0,s}})$ is the projection onto $\ker K_s$, and
$\one_{[2\abs h,\sqrt{4h^2+\delta^2}]}(\abs{h_{q^0,s}})=\one_{[0,\delta]}(\abs{K_s})$ for $\delta\ge0$: the
chiral point $E=0$ of the random-flux operator $K_s$ is mapped to the inner band edges $\abs E=2\abs h$;
\item with $R_s:=(K_s^2+4h^2)^{1/2}$, $\e^{-\ii th_{q^0,s}}=\cos(tR_s)-\ii\,h_{q^0,s}R_s^{-1}\sin(tR_s)$;
\item $(K_s^2)_{ii}=J^2\deg(i)$.  If $\TLat$ is a box in $\Z^2$, $K_s^2$ connects only sites of the same
sublattice at distance $2$: $(K_s^2)_{ik}=J^2s_{ij}s_{jk}$ if $k=i\pm2e_\mu$ (with midpoint $j$), and
$(K_s^2)_{ik}=J^2s_{ij_1}s_{j_1k}(1+W_p)$ if $i,k$ are opposite corners of the plaquette $p$ with the other
corners $j_1,j_2$; here $s_{ij}$ is the sign of the edge joining $i$ and $j$.  Thus the diagonal amplitudes of
$K_s^2$ vanish across $\pi$-flux plaquettes.
\end{enumerate}
\end{enumerate}
\end{proposition}

\begin{proof}
(a) $W_\gamma\Psi_0=\varphi_s(\gamma)\Psi_0$ for all $\gamma$, so $P_{q,\varphi}\Psi_0=0$ for $\varphi\neq\varphi_s$.  Write
$\psi=\sum_n\psi(n)\ket n$ with $(-1)^{\abs n}=\Tpi$.  By Step~1 of the proof of \cref{thm:2dmap},
$P_q\ket{n,s}=0$ if $\Tpi_q\neq\Tpi$ and $P_q\ket{n,s}=2^{-(\abs V-1)/2}\Psi_{n,q,s}=2^{-(\abs V-1)/2}V_{q,s}\ket n$
otherwise; sum over $n$.  Since $O$ commutes with all $P_{q,\varphi}$, and $\e^{-\ii tH_\TLat}$ commutes with them,
\[
  \langle\Psi(t),O\Psi(t)\rangle=\sum_q\langle P_{q,\varphi_s}\Psi(t),OP_{q,\varphi_s}\Psi(t)\rangle
  =2^{1-\abs V}\sum_{\Tpi_q=\Tpi}\langle\e^{-\ii tH_q}\psi,O_{q,s}\e^{-\ii tH_q}\psi\rangle,
\]
with $H_q:=\dGamma(h_{q,s})-h\sum_iq_i$ by \eqref{eq:2dsector}; the constant drops out, and
$\e^{-\ii t\dGamma(a)}=\Gamma(\e^{-\ii ta})$.  The uniform law on a hyperplane has density $2$ on its support.

(b) As $\sigma^x_b\ket{+\cdots+}=\ket{+\cdots+}$, $A_i\Psi_0=\Psi_0$, so $G_i\Psi_0=q^0_i\Psi_0$.  Expanding
$\ket{+\cdots+}=2^{-\abs E/2}\sum_s\ket s_z$, the $W$-eigenspace $\varphi$ receives
$P_{q^0,\varphi}\Psi_0=2^{-\abs E/2}\sum_{s'\in O_\varphi}\ket{n^0,s'}$.  For every $\theta$,
\[
  q^0(\theta)\,\varepsilon_{n^0}(\theta)=\prod_{i\in V_\theta}(-1)^{2n^0_i}=1,
\]
so by \eqref{eq:Psi2d} and \cref{lem:2dgauge}(iii),
$\Psi_{n^0,q^0,s}=2^{-(\abs V-1)/2}\sum_{s'\in O_\varphi}\ket{n^0,s'}$ for every $s\in O_\varphi$.  Hence $P_{q^0,\varphi}\Psi_0=2^{(\abs V-1-\abs E)/2}\Psi_{n^0,q^0,s}=2^{-\beta/2}V_{q^0,s}\ket{n^0}$.
The uniformity statements on boxes and tori follow from \cref{thm:2dmap}(c),(d), since the coordinates of a
uniformly distributed element of $\{\pm1\}^k$ with respect to a basis are i.i.d.\ uniform.  Formula
\eqref{eq:2dlawx} follows as in (a).

(c) Every edge joins opposite sublattices, so $\Gamma_\epsilon K_s\Gamma_\epsilon=-K_s$, and
$q^0_i=(-1)^{n^0_i}=-\epsilon_i$, so $h_{q^0,s}=K_s-2h\Gamma_\epsilon$.  (i)
$(K_s-2h\Gamma_\epsilon)^2=K_s^2-2h(K_s\Gamma_\epsilon+\Gamma_\epsilon K_s)+4h^2\Gamma_\epsilon^2=K_s^2+4h^2$.  (ii) By
(i), $h_{q^0,s}^2\ge4h^2$, which gives the spectral inclusion, and by the functional calculus
$\abs{h_{q^0,s}}=(K_s^2+4h^2)^{1/2}=g(\abs{K_s})$; hence $\one_B(\abs{h_{q^0,s}})=(\one_B\circ g)(\abs{K_s})
=\one_{g^{-1}(B)}(\abs{K_s})$.  Take $B=\{2\abs h\}$ and $B=[2\abs h,\sqrt{4h^2+\delta^2}]$, for which
$g^{-1}(B)=\{0\}$ and $[0,\delta]$.  (iii) For real $E\ne0$, $\e^{-\ii tE}=\cos(t\abs E)-\ii(E/\abs E)\sin(t\abs E)$;
apply the functional calculus to $h_{q^0,s}$, whose spectrum avoids $(-2\abs h,2\abs h)$, and use
$\abs{h_{q^0,s}}=R_s$.  (iv) $(K_s^2)_{ik}=J^2\sum_js_{ij}s_{jk}$, summed over the common neighbours $j$ of $i$
and $k$.  In a box these exist only for $k=i$ ($\deg(i)$ terms equal to $1$), for $k=i\pm2e_\mu$ (one common
neighbour), and for diagonal pairs, which have the two common neighbours $j_1,j_2$; then
$s_{ij_2}s_{j_2k}=W_p\,s_{ij_1}s_{j_1k}$ because $W_p=s_{ij_1}s_{j_1k}s_{kj_2}s_{j_2i}$ and $s^2=1$.
\end{proof}

\begin{remark}[Interpretation]\label{rem:2dlaws}
\leavevmode
\begin{enumerate}[label=(\alph*)]
\item \Cref{prop:2dlaws}(a) is the two-dimensional analogue of \cref{cor:sector}: for $\sigma^z$-polarized
links all gauge-invariant dynamics is an average over the zero-flux Anderson model $h_{q,+}=-J\Delta+2h\diag(q)$
with Bernoulli potential, restricted to a parity hyperplane of charge configurations.  This is the setting of
\cite{SKKM2018}.  A $\sigma^z$-product state with flipped links has the same charge law but a deterministic,
in general non-zero, flux $\varphi_s$; if $\varphi_s\equiv1$, then $s=\theta\cdot+$ for some $\theta$
(\cref{lem:2dgauge}), and by \cref{thm:2dmap}(b) the gauge-invariant dynamics is that of the polarized state with
the matter state $\Gamma(D_\theta)\psi$ (a sign for Fock $\psi$), as in \cref{lem:productlaw}(c).
\item In \cref{prop:2dlaws}(b) the disorder sits entirely in the fluxes, while the potential
$2h(-1)^{n^0_i}$ is fixed by the initial matter configuration.  The state $\ket{+\cdots+}$ has $A_i\equiv1$; the
toric-code ground states considered in \cite[Sect.~I.C]{SKKM2018} are its normalized projections onto single flux
sectors, in which neither the charges nor the fluxes are random.  For the checkerboard (c), the
sector-averaged dynamics is a function of the binary random $\pi$-flux operator $K_s$ alone, by (iii); in
particular every localization question for this initial state is a question about $K_s$ and $K_s^2$.  The
chiral point $E=0$ of $K_s$, where random-flux and random-hopping models are expected to be critical rather than
localized (the Gade singularity \cite{Gade1993}), sits at the inner band edges $\abs E=2\abs h$ of
$h_{q^0,s}$.  We know of no localization theorem for $K_s$ with i.i.d.\ binary fluxes at any energy
(\cref{sec:2d:status}).
\end{enumerate}
\end{remark}

\begin{remark}[Numerical verification]\label{rem:2dED}
The script
\begin{center}
\texttt{code/twod\_check.py}\quad(output: \texttt{code/twod\_check.out})
\end{center}
builds $H_\TLat$, the $G_i$ and the $W_\gamma$ from scratch in the basis $\ket{n,s}$ (Jordan--Wigner fermions, no input from the proof).
\begin{itemize}
\item \emph{\Cref{thm:2dmap}, sector by sector.}  For the planar $2\times3$ box ($\dim\cH=8192$), the
$2\times2$ torus (a multigraph with $\beta=5$: three independent plaquettes and two non-contractible loops;
$\dim\cH=4096$) and the complete graph $K_4$, the joint eigenspaces of all $G_i$ and of the fundamental-cycle
Wilson loops were computed by diagonalizing a random linear combination of them.  Exactly $2^{\abs V+\beta}$
sectors occur, each of dimension $\binom{\abs V}{N_f}$ in the $N_f$-particle space if
$\prod_iq_i=(-1)^{N_f}$ and empty otherwise, and the spectrum of $H_\TLat$ in each of them coincides with that
of $\dGamma(h_{q,s})-h\sum_iq_i$ with $s$ realizing the measured fluxes, to $2\times10^{-14}$ ($J=1$, $h=0.8$).
\item \emph{\Cref{prop:2dlaws}(a),(b).}  For the $2\times4$ and $3\times3$ boxes and the $2\times2$ torus, the
sector weights $\norm{P_{q,\varphi}\Psi_0}^2$ agree with (a) and (b) exactly, and the full time evolution of
$\langle n_x(t)\rangle$ for $t\le8$ ($h\in\{0.4,1.3\}$; checkerboard and other Fock states, both link
polarizations) agrees with \eqref{eq:2dlawz}, \eqref{eq:2dlawx} to $3\times10^{-13}$.
\item \emph{\Cref{prop:2dlaws}(c)} on a $16\times16$ box with random fluxes: (i), (iii) and the entries in (iv) to
$10^{-12}$, and $\min\abs{\sigma(h_{q^0,s})}\ge2\abs h$.
\end{itemize}
\end{remark}

\subsection{Strong coupling: cluster confinement up to times of order \texorpdfstring{$h/J^2$}{h/J²}}\label{sec:2d:strong}

Let $\Lambda\subset\Z^2$ be a box (or $\Lambda=\Z^2$), with nearest-neighbour edges $E(\Lambda)$, and let
$s\in\{\pm1\}^{E(\Lambda)}$, $q\in\{\pm1\}^\Lambda$ be arbitrary; $h_{q,s}=K_s+2h\diag(q)$ acts on $\ell^2(\Lambda)$ as
in \eqref{eq:hqs}.  The \emph{equal-charge clusters} are the connected components of the subgraphs of $\Lambda$
induced on $\{q=+1\}$ and on $\{q=-1\}$; $\Tclus_q(x)$ denotes the cluster containing $x$.  With
$P_\pm:=\one_{\{q=\pm1\}}$ let
\begin{equation}\label{eq:h0}
  h^0_{q,s}:=P_+h_{q,s}P_++P_-h_{q,s}P_-,\qquad h^1_{q,s}:=h_{q,s}-h^0_{q,s}=P_+h_{q,s}P_-+P_-h_{q,s}P_+ ,
\end{equation}
i.e.\ $h^0_{q,s}$ is $h_{q,s}$ with the hopping across all bonds joining opposite charges removed.  For
$x\in\Lambda$ and $R\ge0$ let $B_R(x):=\{y\in\Lambda:\abs{y-x}_\infty\le R\}$.

\begin{proposition}[Cluster confinement at strong coupling]\label{prop:cluster}
Assume $\abs h>2J$ and set $g:=4\abs h-8J>0$.
\begin{enumerate}[label=(\alph*)]
\item $h^0_{q,s}$ leaves $\ell^2(\Tclus)$ invariant for every equal-charge cluster $\Tclus$.
\item For all $t\in\R$, uniformly in $\Lambda$, $q$ and $s$,
\begin{equation}\label{eq:clusterbound}
  \norm{\e^{-\ii th_{q,s}}-\e^{-\ii th^0_{q,s}}}\le\frac{8J}{g}+\frac{48\,\abs t\,J^2}{g} .
\end{equation}
Consequently, for every $x\in\Lambda$, $\norm{\one_{\Lambda\setminus\Tclus_q(x)}\e^{-\ii th_{q,s}}\delta_x}\le8J/g+48\abs tJ^2/g$.
\item Let the law of $q$ have density at most $K$ with respect to the i.i.d.\ Bernoulli$(1/2)$ law on
$\{\pm1\}^\Lambda$ ($K=1$ for i.i.d.\ charges, $K=2$ for the parity hyperplane of \cref{prop:2dlaws}(a)).
There is a numerical constant $c_*>0$ such that for all $x\in\Lambda$, $R\ge4$, $t\in\R$ and $s$,
\begin{equation}\label{eq:clusterprob}
  \E\,\norm{\one_{\Lambda\setminus B_R(x)}\,\e^{-\ii th_{q,s}}\delta_x}^2
  \le\Bigl(\frac{8J}{g}+\frac{48\,\abs t\,J^2}{g}\Bigr)^2+2K\,\e^{-c_*(2R+3)} .
\end{equation}
\end{enumerate}
\end{proposition}

The constant $c_*$ is the one of \cite[Prop.~2.2]{LiL2022}, which rests on $p_c^{\rm site}(\Z^2)>\frac12$
\cite{GrimmettStacey1998} and on the sharpness of the phase transition \cite[Thm.~7.3]{AizenmanBarsky1987};
we quote it in the proof.  No smallness condition on $J/g$ is needed for (b): the bound is proved for all
$\abs h>2J$.  It is informative only when its right-hand side is below $2$, which requires $J/g<\frac14$ and
$\abs t<(2g-8J)/(48J^2)$; in the complementary range (in particular for $J/g\ge\frac14$) it holds trivially,
since the left-hand side never exceeds $2$.  We made no attempt to optimize the constants $8$ and $48$.

\begin{proof}
Replacing $(h,q)$ by $(-h,-q)$ changes neither $h_{q,s}$, $h^0_{q,s}$ nor the clusters, and preserves the laws in
(c); so let $h>0$.  Write $h=h_{q,s}$, $h^0=h^0_{q,s}$, $h^1=h^1_{q,s}$, and let $h^0_\pm:=P_\pm hP_\pm$, regarded as
operators on $\operatorname{Ran}P_\pm$.

(a) The off-diagonal entries of $h^0$ join nearest neighbours of equal charge only; hence
$\langle\delta_y,h^0\delta_x\rangle\ne0$ with $x\in\Tclus$ implies $y\in\Tclus$.

(b) \emph{Norms and gap.}  A self-adjoint matrix with entries of modulus at most $J$ and at most four non-zero
entries in each row has norm at most $4J$ (Schur's test).  This applies to $h^1$ and to $P_\pm K_sP_\pm$, so
$\norm{h^1}\le4J$, $h^0_+=2h+P_+K_sP_+\ge2h-4J$ and $h^0_-=-2h+P_-K_sP_-\le-2h+4J$; the two spectra are separated
by $(2h-4J)-(-2h+4J)=g$.

\emph{Sylvester equation.}  Let $Y:=P_+h^1P_-:\operatorname{Ran}P_-\to\operatorname{Ran}P_+$ and
\[
  X:=\int_0^\infty\e^{-uh^0_+}\,Y\,\e^{uh^0_-}\,\dd u .
\]
The integrand has norm at most $\norm Y\e^{-u(2h-4J)}\e^{u(-2h+4J)}=\norm Y\e^{-ug}$, so the integral converges in
norm and $\norm X\le\norm Y/g\le\norm{h^1}/g$.  Moreover
$h^0_+X-Xh^0_-=-\int_0^\infty\frac{\dd}{\dd u}\bigl(\e^{-uh^0_+}Y\e^{uh^0_-}\bigr)\dd u=Y$.  (This is the solution of
the Sylvester equation for operators with spectra separated by a gap; for general separated spectra
Bhatia--Davis--McIntosh \cite{BDM1983} give the bound $\frac\pi2\norm Y/g$, and the constant $1$ is available
here because $\sigma(h^0_-)$ lies entirely below $\sigma(h^0_+)$.)  Let $S:=X-X^*$, an anti-self-adjoint operator
with $\norm S=\norm X\le\norm{h^1}/g$.  Since $h^1=Y+Y^*$ and
$X^*h^0_+-h^0_-X^*=(h^0_+X-Xh^0_-)^*=Y^*$, we get $[h^0,S]=h^1$.

\emph{Rotation.}  Let $U:=\e^S$, a unitary with $\norm{U-\one}\le\norm S$ (the spectrum of $S$ is imaginary and
$\abs{\e^{\ii\vartheta}-1}\le\abs\vartheta$), and $h':=UhU^*$.  From
$\frac{\dd}{\dd u}\e^{uS}h\e^{-uS}=\e^{uS}[S,h]\e^{-uS}$ and $[S,h]=-[h^0,S]+[S,h^1]=-h^1+[S,h^1]$,
\[
  h'=h^0+\int_0^1\bigl(h^1-\e^{uS}h^1\e^{-uS}\bigr)\dd u+\int_0^1\e^{uS}[S,h^1]\e^{-uS}\dd u .
\]
As $\norm{h^1-\e^{uS}h^1\e^{-uS}}=\norm{\int_0^u\e^{vS}[S,h^1]\e^{-vS}\dd v}\le u\norm{[S,h^1]}$, we obtain
$\norm{h'-h^0}\le\frac32\norm{[S,h^1]}\le3\norm S\norm{h^1}\le3\norm{h^1}^2/g$.  No series is involved.

\emph{Propagators.}  $\e^{-\ii th}=U^*\e^{-\ii th'}U$, and for unitaries $W$,
$\norm{U^*WU-W}\le\norm{U^*W(U-\one)}+\norm{(U^*-\one)W}\le2\norm{U-\one}$.  With Duhamel's formula
$\norm{\e^{-\ii th'}-\e^{-\ii th^0}}\le\abs t\norm{h'-h^0}$,
\[
  \norm{\e^{-\ii th}-\e^{-\ii th^0}}\le\norm{\e^{-\ii th'}-\e^{-\ii th^0}}+\norm{U^*\e^{-\ii th^0}U-\e^{-\ii th^0}}
  \le\frac{3\abs t\norm{h^1}^2}{g}+\frac{2\norm{h^1}}{g},
\]
which gives \eqref{eq:clusterbound} with $\norm{h^1}\le4J$.  By (a), $\e^{-\ii th^0}\delta_x\in\ell^2(\Tclus_q(x))$, so
$\one_{\Lambda\setminus\Tclus_q(x)}\e^{-\ii th}\delta_x=\one_{\Lambda\setminus\Tclus_q(x)}(\e^{-\ii th}-\e^{-\ii th^0})\delta_x$.

(c) Let $\epsilon(t)$ be the right-hand side of \eqref{eq:clusterbound}.  If $\Tclus_q(x)\subset B_R(x)$, then
$\norm{\one_{\Lambda\setminus B_R(x)}\e^{-\ii th}\delta_x}\le\epsilon(t)$ by (b); otherwise the norm is at most $1$.
Hence the left side of \eqref{eq:clusterprob} is at most $\epsilon(t)^2+K\,\Prob_{\rm iid}(\Tclus_q(x)\not\subset B_R(x))$.
Extend $q$ to an i.i.d.\ field on $\Z^2$; the cluster of $x$ in $\Lambda$ is contained in its cluster in $\Z^2$.
If the latter contains a point $y$ with $\abs{y-x}_\infty>R$, a path inside the cluster from $x$ to $y$ changes
$\abs{\cdot-x}_\infty$ by at most one per step, so it reaches the inner boundary
$\partial^-Q_{2R+3}(x)=\{z:\abs{z-x}_\infty=R+1\}$ of the box $Q_{2R+3}(x)=\{z:\abs{z-x}_\infty\le R+1\}$
along sites of charge $q_x$.  For $q_x=-1$ this is the event $E^{2R+3}_{\rm per}(x)$ of \cite[Prop.~2.2]{LiL2022}, for the i.i.d.\
field $V(a):=\bar V\,\frac{1+q_a}2\in\{0,\bar V\}$ with $\Prob[V(a)=0]=\Prob[V(a)=\bar V]=\frac12$.  That proposition
states that there is a numerical constant $c_*>0$ (called $c_0$ there) with $\Prob[E^l_{\rm per}(b)]<\e^{-c_*l}$ for
every $l>10$ and $b\in\Z^2$.  For $q_x=+1$ the same bound follows from the symmetry $q\mapsto-q$ of the law.
With $l=2R+3>10$, i.e.\ $R\ge4$, we obtain $\Prob_{\rm iid}(\Tclus_q(x)\not\subset B_R(x))<2\e^{-c_*(2R+3)}$.
\end{proof}

\begin{remark}[Meaning for the gauge theory]\label{rem:clusterLGT}
Combining \cref{prop:cluster}(c) with \cref{prop:2dlaws}(a) (density $K=2$): for $\sigma^z$-polarized links, a
box $\Lambda$ and a Fock state $\ket n$ with occupied set $S_n$,
\[
  \sum_{y:\ \dist(y,S_n)>R}\langle\Psi(t),n_y\Psi(t)\rangle
  \le\abs{S_n}\Bigl[\Bigl(\frac{8J}{g}+\frac{48\abs tJ^2}{g}\Bigr)^2+4\e^{-c_*(2R+3)}\Bigr],
\]
because $\langle\Psi(t),n_y\Psi(t)\rangle=2^{1-\abs V}\sum_{\Tpi_q=\Tpi}\sum_{a\in S_n}\abs{(\e^{-\ii th_{q,+}})_{ya}}^2$ and
$\dist(y,S_n)>R$ implies $y\notin B_R(a)$ for every $a\in S_n$ (here $\dist$ is the $\ell^\infty$-distance).
So every particle stays within distance $R$ of its initial position, up to probability
$O((J/g)^2+(tJ^2/g)^2+\e^{-2c_*R})$, as long as $t\ll g/J^2\sim\abs h/J^2$.  The same holds in every flux
sector, since \cref{prop:cluster} is uniform in $s$.  In terms of SKKM's discussion \cite[Sect.~I.C]{SKKM2018}:
at strong coupling the charges, attached to the fermions, are confined to the finite clusters of their own sign.
\end{remark}

\begin{remark}[Why the argument stops at $t\sim\abs h/J^2$]\label{rem:clusterstop}
\leavevmode
\begin{enumerate}[label=(\alph*)]
\item The leading correction in $h'-h^0$ is $\frac12[S,h^1]=O(J^2/g)$.  It is block diagonal with respect to the
charges, and to leading order in $J/h$ its $(+,+)$ block is $YY^*/(4h)$, whose entries $J^2\sum_js_{ij}s_{jk}/(4h)$
couple two sites $i,k$ of equal charge through a common neighbour $j$ of the opposite charge.  These are pairs at
distance $2$: diagonal next-nearest neighbours, or $k=i\pm2e_\mu$.  On the time scale $g/J^2$ these processes
connect different first-order clusters.  The relevant geometry is then that of equal-charge sites with
nearest- \emph{and} next-nearest-neighbour adjacency (the ``$*$-lattice''), and at density $\frac12$ this
percolates: at $p=\frac12$, infinite $*$-clusters of both signs coexist almost surely \cite{Higuchi1982}; numerically,
$p_c^{*}=1-p_c^{\rm site}(\Z^2)\approx0.4073<\frac12$ (matching-lattice duality, \cite{Kesten1982}), while
$p_c^{\rm site}(\Z^2)\approx0.5927>\frac12$.  The classical-percolation mechanism behind \cref{prop:cluster} is therefore
exhausted at $t\sim g/J^2$.  What remains is a quantum-percolation problem of resonances between cluster
eigenstates (the effective couplings $J^2/g$ are much smaller than the intra-cluster level spacings, which are of
order $J$).  Whether it localizes is open; \cite{SKKM2018} discuss a possible quantum-percolation delocalization
transition.
\item On the triangular lattice ($p_c^{\rm site}=\frac12$, critical) and on $\Z^3$ ($p_c^{\rm site}\approx0.31<\frac12$),
the equal-charge clusters at density $\frac12$ are not exponentially small, and even the first step fails.
\end{enumerate}
\end{remark}

\begin{remark}[The checkerboard state at strong coupling]\label{rem:checkerstrong}
For the initial state of \cref{prop:2dlaws}(c) on a box, $q^0=-\epsilon$ alternates, so every equal-charge cluster
is a single site, $h^0_{q^0,s}=2h\diag(q^0)$, and \cref{prop:cluster}(b) gives
$\abs{\langle\delta_x,\e^{-\ii th_{q^0,s}}\delta_x\rangle}\ge1-8J/g-48\abs tJ^2/g$ deterministically, in every
flux sector.  Beyond $t\sim\abs h/J^2$ the dynamics is generated, by \cref{prop:2dlaws}(c)(iii), by
$R_s=(K_s^2+4h^2)^{1/2}$.  Since $0\le1+\frac y2-\sqrt{1+y}\le\frac{y^2}8$ for $y\ge0$ and $\norm{K_s}\le4J$,
\[
  \Bigl\|R_s-2\abs h-\frac{K_s^2}{4\abs h}\Bigr\|\le\frac{\norm{K_s}^4}{64\abs h^3}\le\frac{4J^4}{\abs h^3} .
\]
Hence, up to errors $O(\abs tJ^4/\abs h^3)$, the slow dynamics is that of $K_s^2/(4\abs h)$.  By
\cref{prop:2dlaws}(c)(iv) this is a hopping operator on each sublattice, a square lattice rotated by $45^\circ$
with spacing $\sqrt2$.  Its bonds of length $2$ always carry amplitudes $\pm J^2/(4\abs h)$, while its diagonal
bonds carry $\pm J^2/(2\abs h)$ if $W_p=+1$ and vanish if $W_p=-1$.  So the effective problem is a random-sign
hopping model in which the diagonal bonds are present independently with probability $\frac12$; the on-site terms
are constant in the bulk.  We do not know whether it localizes.
\end{remark}

\subsection{Status of the two-dimensional one-body problems}\label{sec:2d:status}

The analogue of \cref{thm:EC} in two dimensions would be a sector-averaged eigenfunction-correlator bound for
$h_{q,+}$ (Bernoulli potential, zero flux) or for $K_s$ (binary random flux).  Neither is known.  We list what is
available, in the normalization of \eqref{eq:hqs}.  On $\Z^2$ with $q$ i.i.d.\ and $s\equiv+1$, writing
$\Delta_{\rm Lap}:=\Delta-4$ for the graph Laplacian of $\Z^2$,
\begin{equation}\label{eq:2dnormalization}
  h_{q,+}=J\bigl(-\Delta_{\rm Lap}+\bar V\,v_q\bigr)-4J-2\abs h,\qquad v_q(x):=\tfrac12\bigl(1+\operatorname{sgn}(h)\,q_x\bigr)\in\{0,1\},\quad
  \bar V:=\frac{4\abs h}J=2\abs\lambda ,
\end{equation}
so $h_{q,+}$ is, up to scale and shift, the Anderson--Bernoulli model $-\Delta_{\rm Lap}+\bar Vv$ of
\cite{DingSmart2020,LiL2022}, with i.i.d.\ fair $v_x\in\{0,1\}$.  (Conjugation by $(-1)^{x_1+x_2}$ maps $\Delta$ to
$-\Delta$ and, combined with $q\mapsto-q$, maps $h_{q,+}$ to $-h_{-q,+}$, so the law of the spectrum is symmetric.)

\begin{enumerate}[label=(\arabic*)]
\item \emph{Band edges} (Ding--Smart \cite[Thm.~1.1]{DingSmart2020}).  For $d=2$ and every $\delta>0$ there is
$\varepsilon>0$ such that, almost surely, $-\Delta_{\rm Lap}+\delta v$ has Anderson localization in $[0,\varepsilon]$:
every polynomially bounded generalized eigenfunction with energy in $[0,\varepsilon]$ decays exponentially.  By
\eqref{eq:2dnormalization} and the symmetry above, almost surely $h_{q,+}$ has Anderson localization in windows of
width $J\varepsilon(2\abs\lambda)$ at both edges $\pm(4J+2\abs h)$ of its spectrum.  This is a spectral statement for
a part of the spectrum that carries little weight of the spectral measures of the vectors $\delta_x$; it gives no
dynamical bound for Fock initial states.
\item \emph{Large disorder} (L.~Li \cite[Thm.~1.1]{LiL2022}).  For $d=2$ and $p=\frac12$ there are finitely many
energies $\lambda^{(1)},\dots,\lambda^{(n)}\in[0,8]$, Dirichlet eigenvalues of $-\Delta_{\rm Lap}$ restricted to finite
connected subsets of $\Z^2$, such that for all sufficiently large $\bar V$, almost surely $-\Delta_{\rm Lap}+\bar Vv$
has Anderson localization in its spectrum outside the windows $[\lambda^{(i)}-\bar V^{-1/4},\lambda^{(i)}+\bar V^{-1/4}]$
and their images under $x\mapsto\bar V+8-x$.  The excluded energies are of the same nature as the energies of the
cluster-confined states of \cref{prop:cluster}: for an equal-charge cluster $\Tclus$ of sites with $v=0$,
$h^0_{q,+}|_\Tclus=J(-\Delta_{\rm Lap})^{\rm Dir}_\Tclus-4J-2\abs h$, so in the variables of \cite{LiL2022} the cluster
energies are Dirichlet eigenvalues of finite connected sets (those of the sets used in \cite[Def.~2.12]{LiL2022}
are excluded).  These windows carry a non-vanishing part of the local spectral weight.  For example, let $x$ be a site whose four neighbours have the opposite charge (probability
$\frac1{16}$), so that $\Tclus_q(x)=\{x\}$, and let $q_x=-\operatorname{sgn}h$.  With $U$, $h'$ as in the proof of
\cref{prop:cluster} and $I:=[-2\abs h-\delta,-2\abs h+\delta]$, the spectral measure $\mu_x$ of $\delta_x$ for $h_{q,+}$
satisfies
\[
  \mu_x(\R\setminus I)^{1/2}=\norm{\one_{\R\setminus I}(h')U\delta_x}\le\norm{U-\one}+\frac{\norm{(h'+2\abs h)\delta_x}}\delta
  \le\frac{4J}g+\frac{48J^2}{g\delta},
\]
since $h^0\delta_x=-2\abs h\,\delta_x$ and $\norm{h'-h^0}\le48J^2/g$.  For Li's window,
$\delta=J\bar V^{-1/4}$, the right side is $O((J/\abs h)^{3/4})$.  So, at large $\abs\lambda$, almost all of
$\mu_x$ lies in an excluded window, and \cite{LiL2022} gives no information about the dynamics started at $x$.
\item \emph{Binary random $\pi$-flux.}  For $K_s$ with i.i.d.\ uniform plaquette fluxes $W_p\in\{\pm1\}$
(\cref{prop:2dlaws}(b),(c)), we know of no localization result at any energy.  The Wegner-type estimates of
\cite{DingSmart2020,LiL2022} rely on eigenvalue variation under changes of the Bernoulli potential at ``free''
sites; flipping the sign of a hopping amplitude is not a monotone perturbation, and we are not aware of an
extension to this case.  By
\cref{prop:2dlaws}(c)(ii), the chiral point of $K_s$ is exactly what matters for the checkerboard initial state
near $\abs E=2\abs h$.
\item \emph{No compactly supported eigenstates} (\cref{prop:nocompact} below).  On the square lattice, strict
localization of an eigenstate in a finite region is impossible for every potential and every flux pattern.  So,
unlike Aharonov--Bohm caging on the dice lattice at flux $\pi$ \cite{VidalMosseriDoucot1998}, no flux sector can
produce a trivially localized (flat-band) sector-dependent theorem.
\end{enumerate}

\begin{proposition}[No compactly supported eigenvectors on $\Z^2$]\label{prop:nocompact}
Let $a_{xy}\in\C\setminus\{0\}$ for every nearest-neighbour pair $\{x,y\}$ of $\Z^2$, let $V:\Z^2\to\C$, and let
$(Hu)(x):=\sum_{y\sim x}a_{xy}u(y)+V(x)u(x)$ for finitely supported $u$.  If $u$ is finitely supported and $Hu=Eu$ for
some $E\in\C$, then $u=0$.  In particular, for no $q$ and no $s$ does $h_{q,s}$ on $\ell^2(\Z^2)$ have a finitely
supported eigenvector.
\end{proposition}

\begin{proof}
Suppose $u\neq0$ and let $S:=\supp u$, a finite set.  Let $m:=\max\{z_1+z_2:z\in S\}$, and among the points $z\in S$
with $z_1+z_2=m$ choose the one with the largest $z_1$.  Let $y:=z+e_1\notin S$.  Its neighbours are $y+e_1$ and
$y+e_2$ (coordinate sum $m+2$), $y-e_2=z+e_1-e_2$ (coordinate sum $m$ and first coordinate $z_1+1$), and $y-e_1=z$;
only $z$ lies in $S$.  The eigenvalue equation at $y$ reads $0=Eu(y)=(Hu)(y)=a_{yz}u(z)\neq0$, a contradiction.
\end{proof}

\begin{remark}[Majorana matter]\label{rem:majorana}
If the complex fermions are replaced by Majorana fermions hopping in a static $\Z_2$ gauge field, as in Kitaev's
honeycomb model \cite{Kitaev2006}, then random flux configurations put the matter in the symmetry class D.  There
numerical studies find a delocalized ``thermal metal'' \cite{Laumann2012,Self2019}.  So the two-dimensional
Majorana analogue of the disorder-free localization discussed here is expected to \emph{fail}.  This is a
statement about numerics and field theory, not a theorem; we mention it only to stress that binary gauge disorder
in two dimensions does not localize automatically.
\end{remark}
